%!TEX root = ../../main.tex
%% ===============================================================
%% 4.8 검출기 타당성 검증 설계
%% 파일: chapters/04_localization/08_validation_design.tex
%% 상위: chapters/04_localization/chapter.tex
%% 정의 라벨: sec:loc-validation
%% 참조 라벨: alg:localization, sec:results-detector
%% 포함 객체: -
%% 인용: -
%% 상태: 초안 (docs/OUTLINE.md 참고)
%% ===============================================================

\section{검출기 타당성 검증 설계}
\label{sec:loc-validation}

알고리즘 \ref{alg:localization}은 사람의 주석 없이 말뭉치를 만들지만, 그 결과가 사람이 ``교전''이라 부르는 것과 일치하는지는 별도로 검증해야 한다(RQ1). 본 논문은 두 트랙의 검증을 설계하였다.

\paragraph{트랙 A: 텔레메트리 재생 기반 주석 연구.} 검출기가 소비하는 것과 동일한 Match-V5 텔레메트리로 미니맵 재생을 재구성하고, 주석자가 검출기 출력을 보지 않은 상태(blind)에서 모든 교전 구간을 표시한다. 이 트랙은 알고리즘의 \emph{결정 규칙}(18 초, 4{,}000 단위, 1{,}800 단위, 15 초/2{,}000 단위)이 사람의 판단과 일치하는지를 검증한다. 다만 API에 없는 정보(킬 없는 대치)는 검출할 수 없으므로 별도 항목(\texttt{killless\_annotated})으로 보고한다.

\paragraph{트랙 B: 전체 리플레이 골드셋.} 2026년 KR 랭크 경기(공개 패치 \goldPatchPublic{}, API 패치 \goldPatchApi{}) \numReplays{} 건의 완전한 리플레이(\texttt{.rofl})를 수집하고, 리플레이·detail·timeline 세 파일의 SHA-256을 대조하였다. 이 가운데 10 분 미만 경기 \numReplaysExcluded{} 건을 제외하고, 생성 시각 \numGoldBins{} 개 구간에서 각 10 경기씩 결정론적으로(시드 \goldSeed{}) \numGoldMatches{} 경기의 골드셋을 뽑았다. 표본 추출에 검출기 출력은 사용하지 않았다. 검출기는 이 \numGoldMatches{} 경기에서 \numGoldCandidates{} 개의 후보 구간(경기당 \goldCandidatesPerMatch{} 개)을 산출하였고, 실행 예외·반복 실행 불일치·경계 오류는 없었다.

두 트랙 모두 두 명의 주석자가 독립적으로 교전의 \emph{시작}(첫 킬이 아닌 교전 개시)부터 \emph{종료}(이탈 또는 전멸)까지를 표시하며, 결과는 시간 IoU 0.5 기준의 탐욕적 일대일 짝짓기로 정밀도·재현율·$F_1$·평균 IoU·평균 온셋 오차를 계산하고, 주석자 간 $F_1$과 Cohen의 $\kappa$를 함께 보고한다. 트랙 B의 사람 주석은 본 논문 작성 시점에 진행 중이며(제 \ref{sec:results-detector} 절), 기술적 통과 자체를 검출 정확도로 인용하지 않는다.
